\chapter{B-Trees}

Up to now, every search structure in this book has relied on a simplifying
assumption: all your data fits in main memory, and reaching any node costs roughly
the same as reaching any other. In RAM, counting each pointer-follow as one unit
of work is a reasonable approximation. Once your dataset outgrows memory and
spills onto secondary storage -- a magnetic disk or a solid-state drive -- that
uniform-cost assumption breaks down. Suddenly \emph{where} a node lives matters enormously, and the cost of a
search is dominated not by how many comparisons you do but by how many times you
have to go fetch a chunk of data from far away.

This chapter is about the search tree that takes that cost asymmetry
seriously. A B-tree generalizes a binary search tree by trading height for width:
instead of tall, narrow two-way branching, it uses wide, shallow multi-way nodes. Instead of
one key and two children per node, a B-tree node holds dozens or hundreds of keys
and hundreds or thousands of children. The result is a tree so shallow that you
can index billions of records and still reach any one of them in a handful of
disk visits. Most relational databases and many modern filesystems organize
their data with a B-tree or a close relative. We will build
one from the ground up.

\section{Why Secondary Storage Dictates Wide Trees}

The design of the B-tree follows from one physical fact about storage hardware.

\begin{intuitionbox}
Imagine your data is a giant library. Fast memory (RAM) is the desk in front of
you: tiny, but anything on it is instantly available. The disk is the warehouse
across town: enormous, but every trip to fetch a box takes ages compared to
reading something already on your desk. Importantly, the warehouse does not hand you
\emph{one page} per trip -- it hands you a whole \emph{box}, because the trip
itself is the expensive part, not the box. So the smart move is: design your data
so that each warehouse trip brings back a box stuffed with exactly the things you
will need next, and so that you need as few trips as possible.
\end{intuitionbox}

That ``box'' has a name. Disks (and the operating system on top of them) move data
in fixed-size blocks called \emph{pages} -- typically something like 4 kilobytes,
sometimes much larger. You cannot read a single byte from disk; you read a whole
page. And the cost of reading that page is dominated by the \emph{seek}: physically
moving a read head, or waiting out the latency of a flash lookup. Once the page is
on its way, the difference between grabbing 8 bytes from it and grabbing 4000 bytes
from it is negligible.

The numbers make the asymmetry clear.

\begin{aestheticbox}
A modern CPU performs a few billion simple operations per second, so one
comparison costs on the order of a nanosecond ($10^{-9}$ seconds). A single random
read from a spinning disk costs on the order of $10$ milliseconds ($10^{-2}$
seconds). That is a factor of roughly \emph{ten million}. One disk access is worth
about ten million in-memory comparisons. Solid-state drives shrink the gap, but
even there a random read costs tens of microseconds -- still tens of thousands of
comparisons. The lesson is the same: when your data lives on disk, the number you
must minimize is \emph{disk accesses}, and you should be willing to do a
ridiculous amount of in-memory computation to save even one of them.
\end{aestheticbox}

Now reconsider a balanced binary search tree holding $n = 10^9$ keys (a billion).
Its height is about $\log_2(10^9) \approx 30$. If the tree lives on disk and each
node sits on its own page, a single search touches about $30$ different pages, and
in the worst case that is $30$ disk seeks. At $10$ milliseconds each, that is about
a third of a second to find \emph{one} record. A database answering thousands of
queries per second cannot possibly afford this.

\begin{focusbox}
Two-way branching is the bottleneck here. A binary tree spends a
whole expensive disk access just to make \emph{one} comparison and then branch one
of two ways. We have already paid for an entire page; we should extract far more
than one bit of decision from it. The fix is to cram many keys into a single page
so that one disk access lets us branch one of \emph{many} ways. Fewer, fatter nodes
give a shorter tree, and a shorter tree means fewer disk accesses.
\end{focusbox}

Here is the picture of the trade we are making. On the left, the binary tree:
deep, and each level is a separate disk trip. On the right, a B-tree: each node is
a full page packed with keys, so the tree is only a few levels tall.

\begin{center}
\begin{tikzpicture}[every node/.style={font=\small}]
% binary tree, tall
\node at (-3.4,2.0) {\textbf{binary: many thin levels}};
\foreach \y in {0,...,4}{
  \draw[gray!60] (-5.0,1.4-0.32*\y) rectangle (-1.8,1.62-0.32*\y);
}
\node[align=center] at (-3.4,-0.6) {\footnotesize one comparison\\\footnotesize per disk page};
% btree, short fat
\node at (3.2,2.0) {\textbf{B-tree: few fat levels}};
\draw[sourceblue,thick] (1.4,1.2) rectangle (5.0,1.62);
\foreach \x in {1.4,2.0,2.6,3.2,3.8,4.4}{\draw[gray!50] (\x,1.2)--(\x,1.62);}
\draw[sourceblue,thick] (1.4,0.4) rectangle (5.0,0.82);
\foreach \x in {1.4,2.0,2.6,3.2,3.8,4.4}{\draw[gray!50] (\x,0.4)--(\x,0.82);}
\node[align=center] at (3.2,-0.6) {\footnotesize many keys (and branches)\\\footnotesize per disk page};
\end{tikzpicture}
\end{center}

\begin{historybox}
B-trees were invented by Rudolf Bayer and Edward M.~McCreight in 1970 while working at Boeing Scientific Research Laboratories. Boeing's engineering operations required indexing massive flight test datasets and manufacturing components stored on mechanical magnetic disks and drums. At the time, RAM was prohibitively expensive, and mechanical disk seeks took 10 to 20 milliseconds -- a lifetime compared to microsecond CPU cycles. Bayer and McCreight designed B-trees specifically to minimize the number of disk page accesses. (The meaning of ``B'' has never been settled -- Bayer, Boeing, balanced, or broad; Bayer himself remained famously coy.)

The natural early instinct was to store an in-memory balanced binary tree (such as an AVL tree) directly on secondary storage. This was disastrous: traversing a depth-20 binary search tree required 20 sequential disk seeks, each fetching a 4KB disk page only to inspect a single 32-byte node, taking over 200 milliseconds per lookup. Bayer and McCreight realized that once a 4KB disk sector is read into RAM, scanning dozens or hundreds of keys in memory is essentially free. Widening the tree to have branching factor $t \approx 1000$ shrunk tree height to 3 or 4, slashing disk reads by 80\% and making multi-gigabyte databases practical.

B-trees became the standard index structure in modern systems: relational database indices (PostgreSQL, MySQL InnoDB, SQLite, Oracle), document databases, and high-performance filesystems (ext4, NTFS, macOS APFS, and Btrfs). The main trade-off is concurrency overhead. Because insertions and deletions split or merge nodes along an entire root-to-leaf path, concurrent multi-threaded writes require complex latch crabbing (locking parent nodes until child operations finish). On modern solid-state flash drives (SSDs), random in-place page updates lead to severe write amplification (rewriting entire flash blocks for small updates), which is why log-structured merge (LSM) trees often replace B-trees in write-heavy workloads. Developing cache-oblivious B-trees and write-optimized learned index structures remains an active area of database systems research.
\end{historybox}

\section{What a B-Tree Is}

A B-tree is a rooted tree, but unlike a
binary search tree its nodes are wide. Each node stores a sorted run of keys, and
between and around those keys hang child pointers, one more child than there are
keys.

\begin{intuitionbox}
Picture a single B-tree node as a row of $k$ keys laid out in increasing order,
with $k+1$ ``gaps'' -- one before the first key, one between each adjacent pair,
and one after the last key. Each gap holds a pointer to a child subtree, and that
subtree contains exactly the keys that belong in that gap. So if a node has keys
$[20, 50, 80]$, it has four children: one for keys $< 20$, one for keys strictly
between $20$ and $50$, one for keys strictly between $50$ and $80$, and one for
keys $> 80$. The keys in the node are the \emph{signposts} that tell you which gap
-- which child -- to descend into next.
\end{intuitionbox}

Here is a node with three keys and its four children drawn explicitly. Notice how
each child's allowed range is fenced in by the keys on either side.

\begin{center}
\begin{tikzpicture}[every node/.style={font=\small}]
% the node
\draw[sourceblue,thick] (0,2) rectangle (6,2.6);
\foreach \x in {2,4}{\draw[sourceblue!60] (\x,2)--(\x,2.6);}
\node at (1,2.3) {$20$};
\node at (3,2.3) {$50$};
\node at (5,2.3) {$80$};
% pointer dots in the four gaps
\foreach \x/\lab in {0/a,2/b,4/c,6/d}{
  \fill[targetgreen] (\x,2) circle (1.6pt);
}
% children boxes
\foreach \x/\txt in {-0.6/{$<20$}, 2/{$20..50$}, 4/{$50..80$}, 6.6/{$>80$}}{
  \draw[gray!70] (\x-0.8,0.6) rectangle (\x+0.8,1.2);
  \node at (\x,0.9) {\footnotesize \txt};
}
\draw[->,gray] (0,2) -- (-0.6,1.2);
\draw[->,gray] (2,2) -- (2,1.2);
\draw[->,gray] (4,2) -- (4,1.2);
\draw[->,gray] (6,2) -- (6.6,1.2);
\end{tikzpicture}
\end{center}

To keep the tree shallow and balanced we impose limits on how empty or how full a
node may be. The key parameter is a fixed integer $t \ge 2$ called the
\emph{minimum degree} of the tree. It is chosen once, when you build the tree,
usually so that a full node exactly fills one disk page.

\begin{formalbox}
\textbf{Definition (B-tree).} Let $t \ge 2$ be a fixed integer, the \emph{minimum
degree}. A \emph{B-tree} is a rooted tree with the following properties.
\begin{enumerate}
\item Every node $x$ stores a number of keys $x.n$, held in nondecreasing sorted
order $x.\mathit{key}_1 \le x.\mathit{key}_2 \le \cdots \le x.\mathit{key}_{x.n}$,
and a boolean $x.\mathit{leaf}$ saying whether $x$ is a leaf.
\item An internal (non-leaf) node with $x.n$ keys has exactly $x.n + 1$ children,
pointed to by $x.c_1, \ldots, x.c_{x.n+1}$. Leaves have no children.
\item The keys separate the subtrees: if $k_i$ is any key stored in the subtree
rooted at child $x.c_i$, then
\[
k_1 \le x.\mathit{key}_1 \le k_2 \le x.\mathit{key}_2 \le \cdots \le
x.\mathit{key}_{x.n} \le k_{x.n+1}.
\]
\item \textbf{Every leaf is at the same depth}, which equals the tree's height
$h$. (B-trees grow and shrink only at the root, which is how they stay perfectly
height-balanced.)
\item \textbf{Occupancy bounds.} Every node other than the root has at least
$t-1$ keys, and therefore at least $t$ children if it is internal. Every node has
at most $2t-1$ keys, and therefore at most $2t$ children if it is internal. The
root may have as few as one key (or zero, if the tree is empty), but obeys the
same upper bound.
\end{enumerate}
A node with the maximum $2t-1$ keys is called \emph{full}.
\end{formalbox}

That occupancy rule -- between $t-1$ and $2t-1$ keys -- is what keeps the
tree balanced. Why are those the right bounds, and not arbitrary ones?

\begin{intuitionbox}
The upper bound ($2t-1$ keys) is so a node fits in one disk page; you pick $t$ to
make a full node fill the page. The lower bound ($t-1$ keys) is so nodes are never
too \emph{empty}. If we let nodes dwindle to one key each, the tree would drift
back toward being binary and tall, defeating the whole purpose. Forcing every node
at least half-full guarantees the branching factor stays large, which keeps the
height tiny. The two bounds together say: \emph{every node is between half-full and
completely full.} The factor-of-two slack between $t-1$ and $2t-1$ is exactly the
room we need to absorb insertions and deletions cheaply, as you will see.
\end{intuitionbox}

\begin{trapbox}
A frequent confusion: ``minimum degree'' $t$ counts \emph{children}, not keys. A
non-root node has \emph{at least} $t$ children and \emph{at most} $2t$ children;
correspondingly it has between $t-1$ and $2t-1$ keys (always one fewer key than
children, for internal nodes). Do not mix up $t$ (a per-node minimum) with the
\emph{order} of a B-tree, a competing piece of terminology some textbooks use for
the maximum number of children. We will stick exclusively with CLRS's minimum
degree $t$ to avoid the muddle. The smallest legal B-tree has $t=2$: nodes hold
$1$ to $3$ keys and have $2$ to $4$ children. Such a tree is sometimes called a
\emph{2-3-4 tree}, and it is the close cousin of the red-black tree from
Chapter 13.
\end{trapbox}

\begin{example}\label{ch18:ex-small-tree}
Let $t = 2$. Then nodes may hold $1$, $2$, or $3$ keys and (if internal) have $2$,
$3$, or $4$ children. Here is a legal B-tree of height $2$ on the keys
$1,2,\ldots,16$. The root holds two keys; each leaf is at the same depth.

\begin{center}
\begin{tikzpicture}[every node/.style={font=\footnotesize}, level distance=14mm]
\node[draw,sourceblue,thick,inner sep=2pt] (root) {$\;7\;\mid\;13\;$};
\begin{scope}[every node/.style={draw,inner sep=2pt,font=\footnotesize}]
\node (a) at (-4.2,-1.5) {$2\mid4$};
\node (b) at (0,-1.5) {$9\mid11$};
\node (c) at (4.2,-1.5) {$15$};
\end{scope}
\begin{scope}[every node/.style={draw,gray!80,inner sep=1.6pt,font=\scriptsize}]
\node (a1) at (-5.6,-3.0) {$1$};
\node (a2) at (-4.2,-3.0) {$3$};
\node (a3) at (-2.8,-3.0) {$5\mid6$};
\node (b1) at (-1.2,-3.0) {$8$};
\node (b2) at (0,-3.0) {$10$};
\node (b3) at (1.2,-3.0) {$12$};
\node (c1) at (3.2,-3.0) {$14$};
\node (c2) at (5.2,-3.0) {$16$};
\end{scope}
\draw (root) -- (a); \draw (root) -- (b); \draw (root) -- (c);
\draw (a)--(a1); \draw (a)--(a2); \draw (a)--(a3);
\draw (b)--(b1); \draw (b)--(b2); \draw (b)--(b3);
\draw (c)--(c1); \draw (c)--(c2);
\end{tikzpicture}
\end{center}

Read the separators: the root's keys $7$ and $13$ split everything into three
ranges -- ``below 7,'' ``between 7 and 13,'' ``above 13'' -- and the three child
subtrees hold exactly those ranges. Within the leftmost child, its keys $2$ and
$4$ split its range further. Every path from root to a leaf has the same length.
\end{example}

\section{How Tall Can a B-Tree Get?}

The point of a B-tree is that it is short, and we can prove exactly how
short. We want an upper bound on the height $h$ in terms of the number of keys $n$,
and we want that bound to be logarithmic with a \emph{large base} -- the larger the
base, the shorter the tree.

The strategy is the one you always use to bound the height of a balanced tree:
count how few nodes a tree of height $h$ can possibly contain. If even the
\emph{sparsest} legal tree of height $h$ must hold a lot of keys, then a tree
holding only $n$ keys cannot be very tall.

\begin{intuitionbox}
To make a B-tree as tall as possible for a given number of keys, you make every
node as \emph{empty} as the rules allow -- because emptier nodes branch less, so it
takes more levels to hold the same keys. The emptiest legal tree has the root with
a single key (its minimum) and every other node with exactly $t-1$ keys and $t$
children (their minimum). So the number of nodes at least \emph{multiplies by} $t$
at each level down. That geometric growth is what forces the height to be only
about $\log_t n$.
\end{intuitionbox}

Count the nodes level by level in this emptiest tree. The root is $1$ node.
The root has at least $2$ children (one key forces two children). Each of those, and
everything below, has at least $t$ children. So the node count is $1$ at the root,
then $2$, and after that it multiplies by at least $t$ at every step down. Figure~\ref{ch18:fig-height-fanout}
shows that geometric growth; the table beside it records the per-level minimum.

\begin{figure}[H]
\begin{center}
\begin{tikzpicture}[every node/.style={font=\footnotesize}]
% level 0
\node[draw,sourceblue,thick,inner sep=2pt] (r) at (0,0) {root};
% level 1: 2 nodes
\node[draw,inner sep=2pt] (a) at (-2.4,-1.3) {$\cdot$};
\node[draw,inner sep=2pt] (b) at (2.4,-1.3) {$\cdot$};
\draw (r)--(a); \draw (r)--(b);
% level 2: t children each -> sketch a fan
\foreach \x in {-3.4,-2.8,-2.0,-1.4}{\draw (a)--(\x,-2.6);\node[draw,gray!70,inner sep=1pt] at (\x,-2.6){\tiny$\cdot$};}
\foreach \x in {1.4,2.0,2.8,3.4}{\draw (b)--(\x,-2.6);\node[draw,gray!70,inner sep=1pt] at (\x,-2.6){\tiny$\cdot$};}
\node at (-2.4,-2.6) {\tiny$\ldots$};
\node at (2.4,-2.6) {\tiny$\ldots$};
% level labels on the right
\node[anchor=west] at (4.1,0)    {depth $0$: $1$ node};
\node[anchor=west] at (4.1,-1.3) {depth $1$: $\ge 2$};
\node[anchor=west] at (4.1,-2.6) {depth $2$: $\ge 2t$};
% brace showing fan-out factor
\draw[decorate,decoration={brace,amplitude=4pt},quotientorange]
  (-3.5,-2.95) -- (-1.3,-2.95);
\node[quotientorange,anchor=north] at (-2.4,-3.05) {\tiny$\ge t$ children};
\end{tikzpicture}
\end{center}
\caption{The sparsest (tallest) B-tree: the root has $\ge 2$ children, and every
node below fans out to $\ge t$ children, so each level holds at least $t$ times as
many nodes as the level above.}
\label{ch18:fig-height-fanout}
\end{figure}

\begin{center}
\begin{tabular}{c c}
\toprule
depth & at least this many nodes \\
\midrule
$0$ & $1$ \\
$1$ & $2$ \\
$2$ & $2t$ \\
$3$ & $2t^2$ \\
$\vdots$ & $\vdots$ \\
$h$ & $2t^{h-1}$ \\
\bottomrule
\end{tabular}
\end{center}

Every node except the root holds at least $t-1$ keys. The number of non-root nodes
is at least the number of nodes at depths $1$ through $h$, which is at least
\[
2 + 2t + 2t^2 + \cdots + 2t^{h-1} = 2\sum_{i=0}^{h-1} t^i = 2\cdot\frac{t^h - 1}{t-1}.
\]
So the total number of keys satisfies
\[
n \;\ge\; \underbrace{1}_{\text{root}} \;+\; (t-1)\cdot 2\cdot\frac{t^h-1}{t-1}
\;=\; 1 + 2(t^h - 1) \;=\; 2t^h - 1.
\]
Rearranging, $t^h \le \dfrac{n+1}{2}$, and taking $\log_t$ of both sides gives the
result.

\begin{formalbox}
\textbf{Theorem (Height of a B-tree).} If $n \ge 1$, then any B-tree of minimum
degree $t \ge 2$ holding $n$ keys has height
\[
h \;\le\; \log_t \frac{n+1}{2}.
\]
\end{formalbox}

\begin{focusbox}
The height is logarithmic in $n$ with base $t$. The large base is what matters.
Compared to a red-black tree, whose height is $\log_2 n$, the B-tree's height is
$\log_t n = \dfrac{\log_2 n}{\log_2 t}$ -- smaller by a factor of $\log_2 t$. With a
realistic $t$ of, say, $1000$, that factor is about $10$, so the B-tree is roughly
ten times shorter, meaning ten times fewer disk accesses per search.
\end{focusbox}

\begin{example}\label{ch18:ex-height-numbers}
Suppose each disk page holds a node with up to $1000$ children, so $2t = 1000$ and
$t = 500$. How tall is the tree for a billion keys, $n = 10^9$?
\[
h \le \log_{500}\frac{10^9 + 1}{2} \approx \log_{500}(5\times 10^8)
= \frac{\ln(5\times 10^8)}{\ln 500} \approx \frac{20.03}{6.21} \approx 3.2,
\]
so $h \le 3$. Three levels! A tree of height $3$ has its root, one level of
internal nodes, and the leaves. If you keep the root permanently cached in memory
(databases always do -- it is just one page), then a search costs at most
\emph{two} disk accesses to reach any of a billion records. Compare that with the
$\approx 30$ accesses a binary tree would demand. In practice that is the
difference between a usable database and an unusable one.
\end{example}

\section{Searching a B-Tree}

Searching is the simplest operation and a good warm-up, because it does not
change the tree at all -- it just walks down it.

\begin{intuitionbox}
At each node you have a sorted row of keys. Scan along it (or binary-search it) to
find where your target key would sit. Either you land exactly on it -- done -- or
you fall into one of the gaps between two keys, and that gap's child pointer tells
you which subtree to descend into. Repeat until you find the key or hit a leaf with
no further child to follow, at which point the key simply is not present. It is the
same higher-or-lower game as a binary search tree, except at each node you choose
among many doors instead of just two.
\end{intuitionbox}

\begin{algorithm}[H]
\caption{\textsc{B-Tree-Search}$(x, k)$ -- find key $k$ in the subtree rooted at $x$}
\KwIn{a node $x$, a search key $k$}
\KwOut{a pair $(y,i)$ such that $y.\mathit{key}_i = k$, or \textsc{nil} if absent}
$i \gets 1$\;
\While{$i \le x.n$ \textbf{and} $k > x.\mathit{key}_i$}{
  $i \gets i + 1$\;
}
\If{$i \le x.n$ \textbf{and} $k = x.\mathit{key}_i$}{
  \Return $(x, i)$\tcp*{found it in this node}
}
\eIf{$x.\mathit{leaf}$}{
  \Return \textsc{nil}\tcp*{nowhere left to go}
}{
  \textsc{Disk-Read}$(x.c_i)$\tcp*{the only line that costs a disk access}
  \Return \textsc{B-Tree-Search}$(x.c_i, k)$\;
}
\end{algorithm}

The single line that matters for cost is \textsc{Disk-Read}: it is the only place
we pay to fetch a page. We perform exactly one disk read per level of descent, so a
search does at most $h + 1 = O(\log_t n)$ disk reads. The in-memory work per node
is the linear scan over up to $2t-1$ keys, so the CPU time is $O(t)$ per node and
$O(t\,h) = O(t \log_t n)$ in total. We are happily trading a lot of in-memory
comparison work for very few disk reads -- exactly the bargain disks demand.

\begin{labbox}
Trace \textsc{B-Tree-Search} for $k = 12$ on the tree of
Example~\ref{ch18:ex-small-tree}.
\begin{itemize}
\item \emph{At the root} $[7,13]$: start $i=1$. Is $12 > 7$? Yes, advance to $i=2$.
Is $12 > 13$? No, stop. Is $12 = 13$? No. The root is not a leaf, so descend into
$c_2$, the child for the range between $7$ and $13$ -- the node $[9,11]$.
\item \emph{At} $[9,11]$: $12 > 9$, advance; $12 > 11$, advance to $i=3$. Now
$i > x.n = 2$, stop. Not equal to anything here. Not a leaf, so descend into $c_3$,
the child for keys above $11$ -- the leaf $[12]$.
\item \emph{At the leaf} $[12]$: $12 = 12$ at $i=1$. \textbf{Found}, return this
node and index $1$.
\end{itemize}
Three nodes visited, one per level -- exactly the height plus one. Now repeat the
trace for $k = 5$ (you should find it in the leaf $[5,6]$) and for the absent
$k = 10\tfrac12$, i.e. a key between $10$ and $11$ (you should fall off the bottom
and return \textsc{nil}).
\end{labbox}

\section{Splitting a Full Node}

Before we can insert, we need one surgical tool. Inserting a new key always
ultimately squeezes it into some leaf, but a leaf might already be full ($2t-1$
keys) with no room. The cure is \emph{splitting}: take a full node and break it in
two, pushing its middle key up into the parent as a new separator.

\begin{intuitionbox}
A full node has $2t-1$ keys. That is an odd number, so it has a unique middle key,
the one at position $t$. Splitting cleaves the node down that middle key: the
$t-1$ keys to its left become a new left node, the $t-1$ keys to its right become a
new right node, and the middle key itself \emph{floats up} to the parent, slotting
in as a brand-new separator between those two new children. Each of the two halves
has exactly $t-1$ keys -- the legal minimum -- so both are valid nodes, and the
parent gains one key and one child. Think of it as a cell dividing: one full cell
becomes two half-cells plus a signpost handed to the level above.
\end{intuitionbox}

Figure~\ref{ch18:fig-split-general-and-t2} illustrates the geometry of the split operation for general parameter $t$ as well as the special case $t=2$ ($2$-$3$-$4$ tree).

\begin{figure}[h]
\centering
\begin{tikzpicture}[font=\small]
  % Lane 1: General t split
  \begin{scope}
    \node[font=\footnotesize\scshape, text=gray!70!black, anchor=west] at (-1.0, 1.8) {General Parameter $t$: Splitting a Full Node of $2t-1$ Keys};

    % Full node before
    \draw[thick, draw=sourceblue, fill=sourceblue!10, rounded corners=1pt] (0, 0) rectangle (5.0, 0.7);
    \node[font=\tiny] at (1.0, 0.35) {$t-1$ keys};
    \draw[thick, draw=sourceblue] (2.0, 0) -- (2.0, 0.7);
    \node[draw=trapred, fill=trapred!25, thick, rounded corners=1pt, font=\scriptsize\bfseries, inner sep=2pt] at (2.5, 0.35) {$k_t$};
    \draw[thick, draw=sourceblue] (3.0, 0) -- (3.0, 0.7);
    \node[font=\tiny] at (4.0, 0.35) {$t-1$ keys};

    % Arrow
    \draw[-{Latex[length=2.5mm]}, very thick, quotientorange] (5.4, 0.35) -- 
      node[above, font=\tiny, text=quotientorange] {Elevate $k_t$} (6.8, 0.35);

    % After: median promoted, two nodes of t-1 keys
    \begin{scope}[xshift=7.2cm]
      % Parent with elevated kt
      \draw[thick, draw=gray!60, fill=gray!8, rounded corners=1pt] (1.3, 0.9) rectangle (2.7, 1.5);
      \node[draw=trapred, fill=trapred!25, thick, rounded corners=1pt, font=\scriptsize\bfseries, inner sep=2pt] at (2.0, 1.2) {$k_t$};

      % Two children of t-1 keys each
      \draw[thick, draw=sourceblue, fill=sourceblue!10, rounded corners=1pt] (0, -0.4) rectangle (1.8, 0.3);
      \node[font=\tiny] at (0.9, -0.05) {$t-1$ keys};

      \draw[thick, draw=sourceblue, fill=sourceblue!10, rounded corners=1pt] (2.2, -0.4) rectangle (4.0, 0.3);
      \node[font=\tiny] at (3.1, -0.05) {$t-1$ keys};

      \draw[-{Latex[length=1.5mm]}, gray!70] (1.7, 0.9) -- (0.9, 0.3);
      \draw[-{Latex[length=1.5mm]}, gray!70] (2.3, 0.9) -- (3.1, 0.3);
    \end{scope}
  \end{scope}

  % Lane 2: Special case t=2 (2-3-4 tree)
  \begin{scope}[yshift=-2.8cm]
    \node[font=\footnotesize\scshape, text=gray!70!black, anchor=west] at (-1.0, 1.8) {Special Case $t=2$ ($2$-$3$-$4$ Tree): $2t-1 = 3$ Keys};

    % 3-key node before
    \draw[thick, draw=sourceblue, fill=sourceblue!10, rounded corners=1pt] (0.5, 0) rectangle (4.5, 0.7);
    \node[font=\scriptsize] at (1.1, 0.35) {$k_1$};
    \draw[thick, draw=sourceblue] (1.8, 0) -- (1.8, 0.7);
    \node[draw=trapred, fill=trapred!25, thick, rounded corners=1pt, font=\scriptsize\bfseries, inner sep=2pt] at (2.5, 0.35) {$k_2$};
    \draw[thick, draw=sourceblue] (3.2, 0) -- (3.2, 0.7);
    \node[font=\scriptsize] at (3.9, 0.35) {$k_3$};

    % Arrow
    \draw[-{Latex[length=2.5mm]}, very thick, quotientorange] (5.4, 0.35) -- 
      node[above, font=\tiny, text=quotientorange] {Elevate $k_2$} (6.8, 0.35);

    % After
    \begin{scope}[xshift=7.2cm]
      \draw[thick, draw=gray!60, fill=gray!8, rounded corners=1pt] (1.5, 0.9) rectangle (2.5, 1.5);
      \node[draw=trapred, fill=trapred!25, thick, rounded corners=1pt, font=\scriptsize\bfseries, inner sep=2pt] at (2.0, 1.2) {$k_2$};

      \draw[thick, draw=sourceblue, fill=sourceblue!10, rounded corners=1pt] (0.3, -0.4) rectangle (1.5, 0.3);
      \node[font=\scriptsize] at (0.9, -0.05) {$k_1$};

      \draw[thick, draw=sourceblue, fill=sourceblue!10, rounded corners=1pt] (2.5, -0.4) rectangle (3.7, 0.3);
      \node[font=\scriptsize] at (3.1, -0.05) {$k_3$};

      \draw[-{Latex[length=1.5mm]}, gray!70] (1.7, 0.9) -- (0.9, 0.3);
      \draw[-{Latex[length=1.5mm]}, gray!70] (2.3, 0.9) -- (3.1, 0.3);
    \end{scope}
  \end{scope}
\end{tikzpicture}
\caption{B-tree node splitting for general parameter $t$ (top) and the $t=2$ case (bottom). In both configurations, a full node holding $2t-1$ keys splits at its median element $k_t$, which is promoted to the parent. The remaining $2t-2$ keys split symmetrically into two legal minimal nodes of $t-1$ keys each.}
\label{ch18:fig-split-general-and-t2}
\end{figure}

Comparing the parent before and after the split, it gains one key (the
promoted median $k_t$) and one child pointer. For this to be safe, \textbf{the parent must
not itself be full} when we ask it to absorb a new key -- otherwise it would
overflow. That single requirement is what shapes the entire insertion algorithm.
The procedure assumes its caller hands it a non-full parent $x$ and the index $i$
of the full child to split.

\begin{algorithm}[H]
\caption{\textsc{B-Tree-Split-Child}$(x, i)$ -- split the full child $x.c_i$}
\KwIn{a non-full internal node $x$ and an index $i$ with $x.c_i$ full}
\KwOut{$x.c_i$ split in two, its median key moved up into $x$}
$y \gets x.c_i$\tcp*{the full child, with $2t-1$ keys}
$z \gets$ a new node; \quad $z.\mathit{leaf} \gets y.\mathit{leaf}$; \quad $z.n \gets t-1$\;
\For{$j \gets 1$ \KwTo $t-1$}{
  $z.\mathit{key}_j \gets y.\mathit{key}_{j+t}$\tcp*{copy $y$'s upper $t-1$ keys into $z$}
}
\If{\textbf{not} $y.\mathit{leaf}$}{
  \For{$j \gets 1$ \KwTo $t$}{
    $z.c_j \gets y.c_{j+t}$\tcp*{and $y$'s upper $t$ children}
  }
}
$y.n \gets t-1$\tcp*{$y$ keeps only its lower $t-1$ keys}
\For{$j \gets x.n+1$ \KwTo $i+1$ \textbf{(decreasing)}}{
  $x.c_{j+1} \gets x.c_j$\tcp*{shift $x$'s children right to open a slot for $z$}
}
$x.c_{i+1} \gets z$\;
\For{$j \gets x.n$ \KwTo $i$ \textbf{(decreasing)}}{
  $x.\mathit{key}_{j+1} \gets x.\mathit{key}_j$\tcp*{shift $x$'s keys right to open a slot}
}
$x.\mathit{key}_i \gets y.\mathit{key}_t$\tcp*{the median floats up into $x$}
$x.n \gets x.n + 1$\;
\textsc{Disk-Write}$(y)$; \textsc{Disk-Write}$(z)$; \textsc{Disk-Write}$(x)$\;
\end{algorithm}

Every loop runs at most $t$ times, so the whole split is $O(t)$ in-memory work plus
three disk writes -- a tidy $O(1)$ in disk accesses. Splitting never
touches anything below $y$ and $z$; it is purely local surgery at one parent and
its one full child.

\begin{trapbox}
Two off-by-one traps lurk here. First, the median sits at index $t$, \emph{not}
$t-1$ or $t+1$; with $2t-1$ keys indexed $1\ldots 2t-1$, position $t$ has exactly
$t-1$ keys on each side. Second, $z$ receives keys $y.\mathit{key}_{t+1}$ through
$y.\mathit{key}_{2t-1}$ -- that is the loop ``$z.\mathit{key}_j \gets
y.\mathit{key}_{j+t}$ for $j=1\ldots t-1$.'' The lower keys
$y.\mathit{key}_1\ldots y.\mathit{key}_{t-1}$ stay in $y$, and we simply set
$y.n = t-1$ to forget the rest. If you accidentally leave the median in one of the
halves, that half ends up with $t$ keys and the parent never receives its
separator, and the whole structure quietly corrupts.
\end{trapbox}

\section{Insertion in a Single Downward Pass}

Now the clever part. The naive plan would be: walk down to the right leaf, insert
there, and if the leaf overflows, split it and push the median up, possibly
triggering a cascade of splits all the way back up to the root. That works, but it
walks the path \emph{twice} -- once down, once back up -- which on disk means
re-reading pages we already paid for. B-trees use a slicker, \emph{single-pass}
strategy that fixes problems on the way down, so it never has to come back up.

\begin{focusbox}
\textbf{The single-pass insight: split full nodes \emph{before} you enter them.}
As you descend toward the target leaf, whenever you are about to step into a child
that is full, split it first. Splitting it pushes a key up into the current node --
but you guaranteed the current node was non-full before you arrived, so it can
safely absorb that key. By induction, every node you stand on during the descent is
non-full, so every split you ever need to do has a safe place to put its median.
You never overflow, you never recurse upward, and you touch each page on the path
exactly once.
\end{focusbox}

There is one node this trick cannot help with: the root, which has no parent to
absorb a pushed-up key. So the root gets special handling. If the root is full when
we begin, we grow a brand-new empty root above it and split the old root into two
children of the new one (Figure~\ref{ch18:fig-rootsplit}). This is the \emph{only}
way a B-tree gains height, and it explains property 4 of the definition -- since
growth happens at the top, all leaves stay at the same depth automatically.

\begin{figure}[H]
\begin{center}
\begin{tikzpicture}[every node/.style={font=\footnotesize}]
% before: a full root (t=2 so 3 keys)
\node at (-3.6,1.5) {\textbf{full root ($t=2$)}};
\draw[sourceblue,thick] (-5.2,0.6) rectangle (-2.0,1.1);
\foreach \x in {-4.13,-3.06}{\draw[sourceblue!50](\x,0.6)--(\x,1.1);}
\node at (-4.66,0.85){$p$};
\node[trapred] at (-3.6,0.85){$q$};
\node at (-2.53,0.85){$r$};
\node at (-0.9,0.85){\large$\Rightarrow$};
% after: new root with q, two children p and r
\node at (3.0,1.9) {\textbf{new root, tree one taller}};
\draw[sourceblue,thick] (2.7,1.2) rectangle (3.3,1.6);
\node[trapred] at (3.0,1.4){$q$};
\draw[sourceblue,thick] (1.6,0.2) rectangle (2.2,0.6);
\node at (1.9,0.4){$p$};
\draw[sourceblue,thick] (3.8,0.2) rectangle (4.4,0.6);
\node at (4.1,0.4){$r$};
\draw[->,gray] (2.85,1.2)--(1.9,0.6);
\draw[->,gray] (3.15,1.2)--(4.1,0.6);
\end{tikzpicture}
\end{center}
\caption{Growing the tree at the root. A full root splits about its median $q$,
which becomes the sole key of a fresh root sitting one level higher; the two halves
$p$ and $r$ become its children. Because the new level appears above everything,
all leaves still share the same depth.}
\label{ch18:fig-rootsplit}
\end{figure}

\begin{algorithm}[H]
\caption{\textsc{B-Tree-Insert}$(T, k)$ -- insert key $k$ into B-tree $T$}
\KwIn{a B-tree $T$ and a key $k$ to insert}
\KwOut{$T$ updated to contain $k$, all invariants preserved}
$r \gets T.\mathit{root}$\;
\eIf{$r.n = 2t-1$}{
  $s \gets$ a new node; \quad $s.\mathit{leaf} \gets \textsc{false}$; \quad $s.n \gets 0$\;
  $s.c_1 \gets r$; \quad $T.\mathit{root} \gets s$\tcp*{new empty root above the old one}
  \textsc{B-Tree-Split-Child}$(s, 1)$\tcp*{split old root; tree grows one level}
  \textsc{B-Tree-Insert-Nonfull}$(s, k)$\;
}{
  \textsc{B-Tree-Insert-Nonfull}$(r, k)$\;
}
\end{algorithm}

\begin{algorithm}[H]
\caption{\textsc{B-Tree-Insert-Nonfull}$(x, k)$ -- insert $k$ into the non-full node $x$ or below it}
\KwIn{a node $x$ that is \emph{not} full, and a key $k$}
$i \gets x.n$\;
\eIf{$x.\mathit{leaf}$}{
  \While{$i \ge 1$ \textbf{and} $k < x.\mathit{key}_i$}{
    $x.\mathit{key}_{i+1} \gets x.\mathit{key}_i$\tcp*{shift keys right to make room}
    $i \gets i - 1$\;
  }
  $x.\mathit{key}_{i+1} \gets k$; \quad $x.n \gets x.n + 1$\;
  \textsc{Disk-Write}$(x)$\tcp*{drop $k$ into the leaf, done}
}{
  \While{$i \ge 1$ \textbf{and} $k < x.\mathit{key}_i$}{
    $i \gets i - 1$\;
  }
  $i \gets i + 1$\tcp*{$x.c_i$ is the child we must descend into}
  \textsc{Disk-Read}$(x.c_i)$\;
  \If{$x.c_i.n = 2t-1$}{
    \textsc{B-Tree-Split-Child}$(x, i)$\tcp*{full child! split it before descending}
    \If{$k > x.\mathit{key}_i$}{
      $i \gets i + 1$\tcp*{the pushed-up median may shift which half we want}
    }
  }
  \textsc{B-Tree-Insert-Nonfull}$(x.c_i, k)$\;
}
\end{algorithm}

That little ``if $k > x.\mathit{key}_i$ then $i \gets i+1$'' is easy to overlook
but essential. After a split, the median key has just landed at position $i$ in
$x$, creating a new child at $i+1$. If our key $k$ is bigger than that median, it
belongs in the right half, so we must redirect into $x.c_{i+1}$ instead of
$x.c_i$.

\subsection{A Fully Traced Insertion}

Watch the machinery run. Take $t = 2$ (so nodes hold $1$ to $3$ keys, and
full means $3$ keys). We start from this tree and insert the key $4$.

\begin{center}
\begin{tikzpicture}[every node/.style={font=\footnotesize}, level distance=14mm]
\node[draw,sourceblue,thick,inner sep=2pt] (root) {$\;10\;$};
\node[draw,inner sep=2pt] (a) at (-2.3,-1.5) {$1\mid3\mid7$};
\node[draw,inner sep=2pt] (b) at (2.3,-1.5) {$14\mid16$};
\draw (root)--(a); \draw (root)--(b);
\end{tikzpicture}
\end{center}

\begin{labbox}
\textbf{Inserting $4$ into the tree above ($t=2$).}
\begin{enumerate}
\item \emph{Check the root.} The root $[10]$ has $1$ key, which is not $2t-1 = 3$,
so the root is not full. We call \textsc{Insert-Nonfull} on it.
\item \emph{At the root, choose a child.} Scanning, $4 < 10$, so we head for the
left child $c_1 = [1,3,7]$. But that child has $3$ keys -- it is \emph{full}! We
must split it before descending.
\item \emph{Split the full child.} The median of $[1,3,7]$ is $3$ (position
$t = 2$). It floats up into the root, which becomes $[3,10]$. The child splits into
$[1]$ (left) and $[7]$ (right). The tree is now:
\begin{center}
\begin{tikzpicture}[every node/.style={font=\footnotesize}, level distance=12mm]
\node[draw,sourceblue,thick,inner sep=2pt] (root) {$\;3\;\mid\;10\;$};
\node[draw,inner sep=2pt] (a) at (-3,-1.3) {$1$};
\node[draw,inner sep=2pt] (b) at (0,-1.3) {$7$};
\node[draw,inner sep=2pt] (c) at (3,-1.3) {$14\mid16$};
\draw (root)--(a);\draw(root)--(b);\draw(root)--(c);
\end{tikzpicture}
\end{center}
\item \emph{Redirect after the split.} Back in the root, the median $3$ now sits at
position $i=1$. Is $k=4 > 3$? Yes, so bump $i$ to $2$: we descend into $c_2 = [7]$,
the right half.
\item \emph{At the leaf.} Node $[7]$ is a non-full leaf. We slide $4$ into place:
$4 < 7$, so $4$ goes before $7$, giving the leaf $[4,7]$. Disk-write it. Done.
\end{enumerate}
The final tree:
\begin{center}
\begin{tikzpicture}[every node/.style={font=\footnotesize}, level distance=12mm]
\node[draw,sourceblue,thick,inner sep=2pt] (root) {$\;3\;\mid\;10\;$};
\node[draw,inner sep=2pt] (a) at (-3,-1.3) {$1$};
\node[draw,inner sep=2pt] (b) at (0,-1.3) {$4\mid7$};
\node[draw,inner sep=2pt] (c) at (3,-1.3) {$14\mid16$};
\draw (root)--(a);\draw(root)--(b);\draw(root)--(c);
\end{tikzpicture}
\end{center}
Notice we proactively split $[1,3,7]$ even though, in hindsight, it had room for
$4$ after the split. The single-pass discipline does not gamble: it splits any full
node it might pass through, guaranteeing the leaf at the bottom always has room.
\end{labbox}

\begin{walkbox}
Here is something to chew on while you stretch your legs. The single-pass insert
sometimes splits a node that, it turns out, did not strictly \emph{need} splitting
-- the eventual insertion happened elsewhere, or the node had spare room. Is that
wasteful? Now think about the alternative: an insert that descends first and splits
on the way back up only when forced. That avoids ``unnecessary'' splits but must
keep the whole root-to-leaf path of pages in memory and re-walk it. On disk, which
cost dominates -- a few extra splits, or re-reading the path? And does your answer
change if the tree lives entirely in RAM? Ask which resource is actually
scarce in each setting; that question is the main lesson of this chapter.
\end{walkbox}

\begin{complexitybox}
\textbf{B-tree operations, minimum degree $t$, $n$ keys, height $h = O(\log_t n)$.}
\begin{center}
\begin{tabular}{l c c}
\toprule
operation & disk accesses & CPU time \\
\midrule
search & $O(h) = O(\log_t n)$ & $O(t\,h) = O(t \log_t n)$ \\
insert & $O(h) = O(\log_t n)$ & $O(t\,h) = O(t \log_t n)$ \\
delete & $O(h) = O(\log_t n)$ & $O(t\,h) = O(t \log_t n)$ \\
\bottomrule
\end{tabular}
\end{center}
Each operation visits $O(h)$ nodes, one disk access per node, and does $O(t)$
in-memory work per node (a scan plus at most one split or merge). Since
$t\log_t n = t\,\frac{\log_2 n}{\log_2 t}$, increasing $t$ shortens the tree but
adds per-node work; you choose $t$ to make a node fill exactly one page, which
minimizes the thing that actually hurts -- disk accesses.
\end{complexitybox}

\section{Why Insertion Is Correct}

It is worth pausing to convince ourselves the single-pass insert really does
preserve every B-tree property. The argument rests on one clean invariant.

\begin{formalbox}
\textbf{Descent invariant.} Every call \textsc{B-Tree-Insert-Nonfull}$(x,k)$ is
made on a node $x$ that is \emph{not full}, i.e. $x.n \le 2t-2$.
\end{formalbox}

\emph{Why the invariant holds.} \textsc{B-Tree-Insert} establishes it for the
first call: if the original root was full, we split it so the new root $s$ holds a
single key (non-full); if it was not full, we call directly on it. Now suppose we
are at a non-full node $x$ and about to recurse into child $x.c_i$. If $x.c_i$ is
full, we split it first. After the split, $x.c_i$ (and its new sibling) each hold
$t-1$ keys -- well below $2t-1$ -- so whichever half we descend into is non-full.
If $x.c_i$ was not full to begin with, we descend into it unchanged, still
non-full. Either way the next call receives a non-full node. By induction the
invariant holds all the way down.

\emph{Why that makes everything correct.} Because the node we land on is always
non-full, the leaf at the bottom always has room to absorb $k$ without overflowing,
so we never violate the upper bound $2t-1$. Every split we perform takes a full
node to two half-full ($t-1$-key) nodes and lifts one key into a guaranteed-non-full
parent, so no node ever drops below $t-1$ keys or exceeds $2t-1$. The sorted order
(property 3) is maintained because we always insert $k$ into the gap where it
belongs and because a split's median is exactly the correct separator between its
two halves. Finally, height grows only when we split the root, and that adds one
level uniformly above everything, so all leaves stay at equal depth (property 4).
Every property survives. $\qed$

\section{Deletion}

Deletion is the fussiest operation, and it is where most first-time
implementers get tangled. The guiding idea, though, mirrors insertion exactly, so
hold onto that symmetry and the cases will feel less like a thicket.

\begin{intuitionbox}
Insertion's enemy was \emph{overflow}: a node getting too full, cured by splitting.
Deletion's enemy is the opposite, \emph{underflow}: a node dropping below $t-1$
keys, cured by either \emph{borrowing} a key from a sibling or \emph{merging} with a
sibling. And just as insertion split full nodes on the way \emph{down} so the leaf
always had room, deletion will \emph{fatten} thin nodes on the way down so that
whenever we finally remove a key, the node we remove it from has a key to spare. We
maintain the mirror-image invariant: \emph{never step into a node with only the
minimum $t-1$ keys; first top it up to at least $t$.}
\end{intuitionbox}

\begin{focusbox}
\textbf{The deletion descent invariant.} Whenever \textsc{B-Tree-Delete} is called
recursively on a node, that node has at least $t$ keys (one more than the minimum
$t-1$). The root is exempt -- it is allowed to be thin -- which is why deletion can
shrink the tree's height only by emptying the root.
\end{focusbox}

To enforce that invariant we need two repair moves. Suppose we want to descend into
a child $c$ that has only $t-1$ keys (the minimum). Before descending, we fatten it:

\begin{itemize}
\item \textbf{Borrow (rotate).} If an \emph{adjacent sibling} of $c$ has at least
$t$ keys (a key to spare), rotate one key across. The separator key in the parent
that sits between $c$ and that sibling moves \emph{down} into $c$, and the
sibling's extreme key (and a child pointer, if internal) moves \emph{up} to take
the separator's place. Now $c$ has $t$ keys and the sibling still has at least
$t-1$. Everybody is legal.
\item \textbf{Merge.} If \emph{both} adjacent siblings have only $t-1$ keys
(nobody can spare a key), merge $c$ with one of them. Pull the separator key down
from the parent and fuse $c$, that separator, and the sibling into one node. The
merged node has $(t-1) + 1 + (t-1) = 2t-1$ keys -- a perfectly full but legal node
-- and the parent loses one key and one child.
\end{itemize}

Figure~\ref{ch18:fig-borrow} shows a borrow and Figure~\ref{ch18:fig-merge-general-and-t2} a
merge; in both, the parent's separator is the hinge the keys rotate around.

\begin{figure}[H]
\begin{center}
\begin{tikzpicture}[every node/.style={font=\footnotesize}]
% BORROW
\node at (-3.6,2.5) {\textbf{borrow from right sibling}};
\draw[gray!70] (-5.0,1.7) rectangle (-2.6,2.15);
\node at (-4.4,1.92){$a$};\node at (-3.2,1.92){$s$};
\draw[sourceblue,thick] (-6.0,0.7) rectangle (-5.0,1.15);
\node at (-5.5,0.92){$c$};
\draw[sourceblue,thick] (-3.7,0.7) rectangle (-2.1,1.15);
\node at (-3.35,0.92){$g$};\node at (-2.45,0.92){$h$};
\draw[->,gray] (-3.2,1.7)--(-3.2,1.15);

\node at (-1.65,1.45){\large$\Rightarrow$};

\draw[gray!70] (-5.0+4.8,1.7) rectangle (-2.6+4.8,2.15);
\node at (-4.4+4.8,1.92){$a$};\node at (-3.2+4.8,1.92){$g$};
\draw[sourceblue,thick] (-6.0+4.8,0.7) rectangle (-5.0+4.8+0.6,1.15);
\node at (-5.5+4.8-0.1,0.92){$c$};\node at (-5.5+4.8+0.7,0.92){$s$};
\draw[sourceblue,thick] (-3.7+4.8,0.7) rectangle (-2.1+4.8,1.15);
\node at (-2.9+4.8,0.92){$h$};
\end{tikzpicture}
\end{center}
\caption{Borrowing. The thin child $c$ takes the parent separator $s$, and the
donor sibling's smallest key $g$ rises to become the new separator.}
\label{ch18:fig-borrow}
\end{figure}

\begin{figure}[h]
\centering
\begin{tikzpicture}[font=\small]
  % Lane 1: General t merge
  \begin{scope}
    \node[font=\footnotesize\scshape, text=gray!70!black, anchor=west] at (-1.0, 2.2) {General Parameter $t$: Merging Two Minimal Siblings ($t-1$ Keys Each)};

    % Before: Parent with separator s and two siblings of t-1 keys
    \draw[thick, draw=gray!60, fill=gray!8, rounded corners=1pt] (1.2, 1.1) rectangle (3.2, 1.7);
    \node[draw=quotientorange, fill=quotientorange!25, thick, rounded corners=1pt, font=\scriptsize\bfseries, inner sep=2pt] at (2.2, 1.4) {$s$};

    \draw[thick, draw=sourceblue, fill=sourceblue!10, rounded corners=1pt] (0, -0.2) rectangle (1.8, 0.5);
    \node[font=\tiny] at (0.9, 0.15) {$t-1$ keys};

    \draw[thick, draw=sourceblue, fill=sourceblue!10, rounded corners=1pt] (2.6, -0.2) rectangle (4.4, 0.5);
    \node[font=\tiny] at (3.5, 0.15) {$t-1$ keys};

    \draw[-{Latex[length=1.5mm]}, gray!70] (1.7, 1.1) -- (0.9, 0.5);
    \draw[-{Latex[length=1.5mm]}, gray!70] (2.7, 1.1) -- (3.5, 0.5);

    % Arrow
    \draw[-{Latex[length=2.5mm]}, very thick, quotientorange] (4.9, 0.6) -- 
      node[above, font=\tiny, text=quotientorange] {Pull $s$ down} (6.5, 0.6);

    % After: Fused into one full node of 2t-1 keys
    \begin{scope}[xshift=7.0cm]
      \draw[thick, draw=sourceblue, fill=sourceblue!10, rounded corners=1pt] (0, 0.1) rectangle (5.0, 0.8);
      \node[font=\tiny] at (1.0, 0.45) {$t-1$ keys};
      \draw[thick, draw=sourceblue] (2.0, 0.1) -- (2.0, 0.8);
      \node[draw=quotientorange, fill=quotientorange!25, thick, rounded corners=1pt, font=\scriptsize\bfseries, inner sep=2pt] at (2.5, 0.45) {$s$};
      \draw[thick, draw=sourceblue] (3.0, 0.1) -- (3.0, 0.8);
      \node[font=\tiny] at (4.0, 0.45) {$t-1$ keys};

      \node[font=\tiny, text=gray!70!black] at (2.5, -0.3) {$(t-1) + 1 + (t-1) = 2t-1$ keys (full node)};
    \end{scope}
  \end{scope}

  % Lane 2: t=2 merge
  \begin{scope}[yshift=-3.2cm]
    \node[font=\footnotesize\scshape, text=gray!70!black, anchor=west] at (-1.0, 2.2) {Special Case $t=2$ ($2$-$3$-$4$ Tree): $t-1 = 1$ Key Each};

    % Before: Parent with separator s and two 1-key siblings
    \draw[thick, draw=gray!60, fill=gray!8, rounded corners=1pt] (1.5, 1.1) rectangle (2.5, 1.7);
    \node[draw=quotientorange, fill=quotientorange!25, thick, rounded corners=1pt, font=\scriptsize\bfseries, inner sep=2pt] at (2.0, 1.4) {$s$};

    \draw[thick, draw=sourceblue, fill=sourceblue!10, rounded corners=1pt] (0.4, -0.2) rectangle (1.4, 0.5);
    \node[font=\scriptsize] at (0.9, 0.15) {$k_1$};

    \draw[thick, draw=sourceblue, fill=sourceblue!10, rounded corners=1pt] (2.6, -0.2) rectangle (3.6, 0.5);
    \node[font=\scriptsize] at (3.1, 0.15) {$k_3$};

    \draw[-{Latex[length=1.5mm]}, gray!70] (1.8, 1.1) -- (0.9, 0.5);
    \draw[-{Latex[length=1.5mm]}, gray!70] (2.2, 1.1) -- (3.1, 0.5);

    % Arrow
    \draw[-{Latex[length=2.5mm]}, very thick, quotientorange] (4.9, 0.6) -- 
      node[above, font=\tiny, text=quotientorange] {Pull $s$ down} (6.5, 0.6);

    % After: Fused into one 3-key node
    \begin{scope}[xshift=7.0cm]
      \draw[thick, draw=sourceblue, fill=sourceblue!10, rounded corners=1pt] (0.5, 0.1) rectangle (4.5, 0.8);
      \node[font=\scriptsize] at (1.1, 0.45) {$k_1$};
      \draw[thick, draw=sourceblue] (1.8, 0.1) -- (1.8, 0.8);
      \node[draw=quotientorange, fill=quotientorange!25, thick, rounded corners=1pt, font=\scriptsize\bfseries, inner sep=2pt] at (2.5, 0.45) {$s$};
      \draw[thick, draw=sourceblue] (3.2, 0.1) -- (3.2, 0.8);
      \node[font=\scriptsize] at (3.9, 0.45) {$k_3$};

      \node[font=\tiny, text=gray!70!black] at (2.5, -0.3) {$1 + 1 + 1 = 3$ keys};
    \end{scope}
  \end{scope}
\end{tikzpicture}
\caption{B-tree node merging for general parameter $t$ (top) and the $t=2$ case (bottom). When two adjacent siblings both hold the minimum $t-1$ keys, their intervening parent separator $s$ is pulled down to fuse all three components into a single valid node containing $(t-1) + 1 + (t-1) = 2t-1$ keys.}
\label{ch18:fig-merge-general-and-t2}
\label{ch18:fig-merge}
\end{figure}

With those two moves in hand, the deletion itself splits into a few cases depending
on where the key to remove lives and whether it is in a leaf.

\begin{formalbox}
\textbf{Deleting key $k$ from the subtree rooted at $x$ (sketch).} First ensure
$x$ has at least $t$ keys (borrow/merge as above, unless $x$ is the root). Then:
\begin{description}
\item[Case 1 -- $k$ is in $x$ and $x$ is a leaf.] Just remove $k$ from $x$. Since
$x$ had $\ge t$ keys it still has $\ge t-1$. Done.
\item[Case 2 -- $k$ is in $x$ and $x$ is internal.] We cannot leave a hole, so we
replace $k$ with a neighboring key that can fill its slot.
\begin{itemize}
\item[\emph{2a.}] If the child $y$ \emph{left} of $k$ has $\ge t$ keys, find $k$'s
\emph{predecessor} $k'$ (the largest key in $y$'s subtree), recursively delete $k'$
from that subtree, and overwrite $k$ with $k'$.
\item[\emph{2b.}] Symmetrically, if the child $z$ \emph{right} of $k$ has $\ge t$
keys, use $k$'s \emph{successor}.
\item[\emph{2c.}] If both flanking children have only $t-1$ keys, \emph{merge} $k$
and the right child into the left child (a full $2t-1$-key node), then recursively
delete $k$ from that merged child.
\end{itemize}
\item[Case 3 -- $k$ is not in $x$.] Find the child $x.c_i$ whose subtree must
contain $k$. If $x.c_i$ has only $t-1$ keys, fatten it first (borrow from a sibling,
or merge), then recurse into the appropriate child to delete $k$.
\end{description}
\end{formalbox}

\begin{intuitionbox}
Case 2 is the only genuinely new idea, and it is the same trick as deleting from an
ordinary binary search tree: to remove an internal key without leaving a gap,
replace it with its \emph{predecessor} or \emph{successor}, then delete \emph{that}
key from where it lives (always lower down, eventually in a leaf). The predecessor
is the rightmost key in the left subtree; the successor is the leftmost key in the
right subtree. We prefer whichever side has a key to spare so the recursion lands on
a node that is already fat enough -- preserving the invariant. If neither side can
spare a key, we merge them and the problem shrinks to one level down.
\end{intuitionbox}

\begin{labbox}
\textbf{A deletion that triggers a merge ($t = 2$).} Start from this tree and delete
the key $6$.
\begin{center}
\begin{tikzpicture}[every node/.style={font=\footnotesize}, level distance=13mm]
\node[draw,sourceblue,thick,inner sep=2pt] (root) {$\;3\;\mid\;9\;$};
\node[draw,inner sep=2pt] (a) at (-3,-1.4) {$1$};
\node[draw,inner sep=2pt] (b) at (0,-1.4) {$6$};
\node[draw,inner sep=2pt] (c) at (3,-1.4) {$11\mid13$};
\draw (root)--(a);\draw(root)--(b);\draw(root)--(c);
\end{tikzpicture}
\end{center}
\begin{enumerate}
\item \emph{At the root} $[3,9]$: $6$ is not here. It lies between $3$ and $9$, so
it belongs in the middle child $c_2 = [6]$.
\item \emph{Fatten before descending.} Child $[6]$ has only $t-1 = 1$ key -- the
minimum. Can it borrow? Its left sibling $[1]$ has $1$ key (nothing to spare); its
right sibling $[11,13]$ has $2$ keys, but that is across the separator $9$, and to
keep the example pointed, suppose we look left first and, finding no donor,
\emph{merge} $[6]$ with its left sibling $[1]$. We pull the separator $3$ down: the
merged node is $[1,3,6]$, and the root loses a key, becoming $[9]$.
\begin{center}
\begin{tikzpicture}[every node/.style={font=\footnotesize}, level distance=13mm]
\node[draw,sourceblue,thick,inner sep=2pt] (root) {$\;9\;$};
\node[draw,inner sep=2pt] (a) at (-2.4,-1.3) {$1\mid3\mid6$};
\node[draw,inner sep=2pt] (c) at (2.4,-1.3) {$11\mid13$};
\draw (root)--(a);\draw(root)--(c);
\end{tikzpicture}
\end{center}
\item \emph{Delete from the fattened leaf.} Now recurse into $[1,3,6]$, a leaf with
$3$ keys. It is Case 1: simply remove $6$, leaving $[1,3]$. Done.
\end{enumerate}
\begin{center}
\begin{tikzpicture}[every node/.style={font=\footnotesize}, level distance=13mm]
\node[draw,sourceblue,thick,inner sep=2pt] (root) {$\;9\;$};
\node[draw,inner sep=2pt] (a) at (-2.4,-1.3) {$1\mid3$};
\node[draw,inner sep=2pt] (c) at (2.4,-1.3) {$11\mid13$};
\draw (root)--(a);\draw(root)--(c);
\end{tikzpicture}
\end{center}
The merge fixed the underflow \emph{on the way down}, so by the time we removed $6$
the leaf had keys to spare. Had the root dropped to zero keys during such a merge,
we would discard it and make the merged node the new root -- that is the lone way a
B-tree loses height.
\end{labbox}

\begin{trapbox}
The classic deletion bug is forgetting to fatten \emph{before} you descend, and
instead deleting first and patching up underflow afterward. The patch-up-afterward
style does work, but it forces you to walk back up the path -- exactly the
double-pass cost we worked so hard to avoid in insertion -- and it is far easier to
get the cascade of merges wrong. The disciplined recipe is always: \emph{guarantee
the node you are about to enter has $\ge t$ keys, then enter it.} If you maintain
that one invariant religiously, every removal you finally perform leaves a legal
node behind, and you never recurse upward.
\end{trapbox}

Like search and insert, deletion visits one node per level and does $O(t)$
in-memory work (one borrow or merge) at each, so it costs $O(h) = O(\log_t n)$ disk
accesses and $O(t\log_t n)$ CPU time -- the same budget as everything else.

\section{B-Trees in Practice}

\begin{aestheticbox}
You now understand the data structure behind most large-scale data storage. When PostgreSQL, MySQL/InnoDB, Oracle, or SQLite build an
index on a column, that index is a B-tree (often the B$^+$-tree variant, which
stores all real data in the leaves and chains the leaves together for fast range
scans -- a small twist on what you have just learned). When you query
``\texttt{WHERE id = 91823}'' and it returns in a millisecond from a table of a
billion rows, you are watching a height-$3$ B-tree do its two disk reads.
Filesystems use them to map filenames to locations. Key-value stores, time-series
databases, and even some in-memory systems use them because the same width-over-height
logic that beats disk latency also beats CPU-cache-miss latency: a cache line is just
a very small, very fast ``page,'' and packing many keys into one
fetched block pays off there too.
\end{aestheticbox}

\begin{focusbox}
The one idea to carry away: \textbf{a B-tree is a balanced search tree re-tuned for a
world where each access is expensive, so it maximizes the work done per access by
making nodes wide.} Width $\Rightarrow$ low height $\Rightarrow$ few accesses. Every
detail -- the $t-1$-to-$2t-1$ occupancy band, splitting full nodes on the way down,
fattening thin nodes on the way down, growing and shrinking only at the root -- serves
that one goal.
\end{focusbox}

B-trees solve the dynamic dictionary problem at large scale, balancing multi-way branching
against the physical page transfers of secondary storage. Every query navigates an ordered key
space to search, insert, or delete individual records.

Many algorithmic problems, however, require managing collections of items where order
does not matter at all. Instead, we need to track dynamic partitions: maintaining disjoint pools of
connected elements, merging pools together, and querying whether two items belong to the same
component. In \textbf{Chapter~19 (Disjoint Sets and Union-Find)}, we study a compact
forest data structure that implements these set operations so efficiently that any sequence of $m$
operations runs in practically linear time -- bounded by the inverse of the fast-growing
Ackermann function.

\section*{Exercises}
\addcontentsline{toc}{section}{Exercises}

\begin{exercise}\label{ch18:ex-warmup-bounds}
For a B-tree with minimum degree $t = 4$, state: (a) the minimum and maximum number
of keys in a non-root node; (b) the minimum and maximum number of children of an
internal non-root node; (c) the maximum number of keys in any node. Then state how
many keys the \emph{root} is allowed to have at minimum.
\end{exercise}

\begin{exercise}\label{ch18:ex-height-count}
Using the height theorem, compute an upper bound on the height of a B-tree with
$t = 100$ holding $n = 10^6$ keys, then $n = 10^{12}$ keys. Roughly how many disk
accesses does a search cost in each case if the root is kept cached in memory?
\end{exercise}

\begin{exercise}\label{ch18:ex-build}
Starting from an empty B-tree with $t = 2$, insert the keys
$F$, $S$, $Q$, $K$, $C$, $L$, $H$, $T$, $V$, $W$, $M$, $R$, $N$, $P$ in this
order (compare them alphabetically). Draw the tree after every insertion that causes a split. (This is
the classic CLRS build-up exercise; doing it by hand cements the single-pass
insert.)
\end{exercise}

\begin{exercise}\label{ch18:ex-minkeys}
Show that a B-tree of minimum degree $t$ and height $h$ holding $n$ keys must
satisfy $n \le (2t)^{h+1} - 1$ as well (the \emph{maximum} number of keys, from the
fullest possible tree), and combine it with the chapter's lower bound to sandwich
$n$. What does the upper bound tell you about the \emph{smallest} possible height
for $n$ keys?
\end{exercise}

\begin{exercise}\label{ch18:ex-split-trace}
Let $t = 3$. First check yourself: how many keys does a full node hold? Write that
full node's keys as $10, 20, 30, 40, 50$, identify the median that floats up under
\textsc{B-Tree-Split-Child}, and write out the two resulting child nodes and the key
handed to the parent.
\end{exercise}

\begin{exercise}\label{ch18:ex-why-not-full-parent}
\textsc{B-Tree-Split-Child} requires that the parent be non-full before it is
called. Explain concretely what goes wrong if you split a child whose parent is
already full, and explain how \textsc{B-Tree-Insert-Nonfull}'s descent guarantees
the parent is never full at the moment of a split.
\end{exercise}

\begin{exercise}\label{ch18:ex-delete-trace}
Take the final tree from the deletion lab in this chapter (the one with root $[9]$,
left child $[1,3]$, right child $[11,13]$, with $t=2$). Delete the key $9$. Which
deletion case applies, and what does the tree look like afterward? (Hint: $9$ lives
in an internal node, so you are in Case 2 -- decide whether to use its predecessor
or successor based on which flanking child can spare a key.)
\end{exercise}

\begin{exercise}\label{ch18:ex-onepass-vs-twopass}
Argue carefully why the single-pass insertion never needs to revisit a node it has
already left on its way down. Phrase your answer in terms of the descent invariant.
Then describe a scenario in which a \emph{two-pass} (descend, then split upward)
insertion would re-read the same disk page multiple times, and count how many extra
disk reads it incurs in the worst case for a tree of height $h$.
\end{exercise}

\begin{exercise}\label{ch18:ex-choose-t}
A disk page holds $4096$ bytes. Each key plus its associated child pointer occupies
$32$ bytes, and a node needs an extra $8$-byte header. Roughly what is the largest
minimum degree $t$ such that a full node still fits in one page? Show your
arithmetic, and explain why you would choose $t$ this large rather than, say,
$t = 2$.
\end{exercise}

\begin{exercise}\label{ch18:ex-bplus}
\emph{(Reading ahead.)} Look up the B$^+$-tree variant. In a B$^+$-tree, all real
records live in the leaves and internal nodes hold only separator keys, and the
leaves are linked in a chain. Explain in one paragraph why this layout makes
\emph{range queries} (``give me every key between $40$ and $90$ in sorted order'')
faster than in the plain B-tree of this chapter.
\end{exercise}
